%==============================================================================
%  CAPÍTULO X
%==============================================================================
\chapter{La paradoja del \texorpdfstring{\emph{discharge}}{discharge} y la exclusión del
Crédito con Aval del Estado}
\label{cap:discharge-cae}

\fichaCapitulo{Incompatibilidad de las deudas de estudio con el perdón concursal.}{Tensión
entre universalidad y especialidad, y evolución de la jurisprudencia de la Corte Suprema.}

%==============================================================================
\bloqueTeorico
%==============================================================================

\subsection{El conflicto de principios}

Colisión entre el beneficio general del \discharge{} (\art{255}) ---dirigido a liberar de
cargas perpetuas al individuo y reinsertarlo productivamente--- y la protección del
interés fiscal comprometido en el financiamiento de la educación superior (\Lcae{}, \cite{ley-20027}; y \Lfondosolidario{},
\cite{ley-19287}).

\subsection{La ausencia de un catálogo legal de deudas excluidas}

El derecho chileno adolece de una carencia que está en la raíz de todo el problema: la
\LIR{} \citeyearpar{ley-20720} no contiene un catálogo de deudas excluidas de la extinción.
El \art{255} declara extinguidos «los saldos insolutos de las obligaciones contraídas por el
deudor con anterioridad al inicio del procedimiento», sin distinguir por naturaleza.

Otros ordenamientos resuelven la tensión en la propia ley: enumeran las deudas que
sobreviven al \discharge{} ---típicamente las alimentarias, las tributarias, las derivadas de
delito y, en varios sistemas, las educativas---. El legislador chileno no lo hizo. La
consecuencia es que la determinación de si una deuda concreta sobrevive quedó entregada a los
tribunales, que han debido construir la respuesta a partir de principios generales, con la
inestabilidad jurisprudencial que documenta la sección \ref{sec:cap10-juris}.

%==============================================================================
\bloqueNormativo
%==============================================================================

\subsection{El efecto extintivo del artículo 255}

El \art{255} es la norma central. Firme la resolución de término, se entienden extinguidos
\destacar{por el solo ministerio de la ley y para todos los efectos legales} los saldos
insolutos de las obligaciones anteriores al inicio del procedimiento, y el deudor se entiende
\destacar{rehabilitado para todos los efectos legales}, salvo que la resolución establezca
algo distinto.

Tres notas de esta redacción condicionan todo el capítulo:

\begin{enumerate}
  \item La extinción opera \emph{de pleno derecho}: no requiere declaración adicional ni
  gestión del deudor.
  \item Alcanza a los \emph{saldos insolutos}: lo pagado en el procedimiento queda firme; lo
  que resta se extingue.
  \item Comprende las obligaciones \emph{anteriores} al inicio: las posteriores no se ven
  afectadas.
\end{enumerate}

\subsection{El requisito previo del artículo 254: no hay término con incidentes pendientes}

El \art{254} contiene una regla de secuencia que la práctica no siempre advierte: si se
promovió el incidente de mala fe del \art{169 A} (capítulo \ref{cap:mala-fe}) o se dedujeron
acciones revocatorias del Capítulo VI (capítulo \ref{cap:revocatorias}), \destacar{el tribunal
no puede dictar la resolución de término} hasta que quede firme la resolución que falla el
incidente o la sentencia sobre las acciones.

\begin{foco}[La extinción está condicionada a la probidad]
El \discharge{} no es automático mientras penda un cuestionamiento a la conducta del deudor.
El art. 254 encadena la extinción de deudas al resultado del incidente de mala fe y de las
revocatorias: primero se depura la conducta, después se extingue la deuda. Es la bisagra que
conecta este capítulo con los capítulos \ref{cap:revocatorias} y \ref{cap:mala-fe}.
\end{foco}

\subsection{Universalidad subjetiva y el principio de especialidad}

El conflicto del CAE se plantea como una colisión entre dos principios. La
\destacar{universalidad subjetiva} del concurso ---que somete a todos los acreedores y todas
las deudas a sus resultados--- empuja hacia la extinción del crédito educativo. La
\destacar{especialidad} de la legislación de educación superior ---\Lcae{}
\citeyearpar{ley-20027} y \Lfondosolidario{} \citeyearpar{ley-19287}--- se invoca para
sustraerlo.

El texto que sostiene esta segunda posición es el \art{8}: «las normas contenidas en leyes
especiales prevalecerán sobre las disposiciones de esta ley». Si la legislación del CAE es
«especial» en el sentido del \art{8}, prevalecería sobre la extinción del \art{255}. La postura
contraria replica con la regla de la \lat{lex posterior} ---la \LIR{}, más reciente, expresaría
la voluntad legislativa última--- y niega que el CAE constituya un «estatuto concursal especial».
La tensión entre ambas lecturas del \art{8} es, en rigor, el núcleo dogmático de todo el
capítulo.

\subsection{El régimen del Crédito con Aval del Estado}

Naturaleza del préstamo, garantías fiscales ejecutables por la \tgr{} y cobro contingente
al ingreso, conforme a la \Lcae{}.

\subsection{El Crédito Solidario Universitario}

Régimen de la \Lfondosolidario{}.

%==============================================================================
\bloquePractico
%==============================================================================

\subsection{Secuencia procesal del deudor con deudas educativas}

\begin{flujograma}{Suerte del CAE en el procedimiento de liquidación voluntaria}{cap10-cae}
  \node[hito] (sol) {EL DEUDOR PRESENTA SOLICITUD DE LIQUIDACIÓN VOLUNTARIA\\[2pt]
    \normalfont Debe declarar la deuda CAE en la nómina de pasivos};
  \node[nodo, below=of sol] (excl)
    {LA INSTITUCIÓN BANCARIA O LA \tgr{} SOLICITA LA EXCLUSIÓN DEL CRÉDITO};

  \coordinate (split) at ($(excl.south)+(0,-6mm)$);
  \node[adverso, below left=12mm and 4mm of excl.south] (acoge)
    {EL TRIBUNAL ACOGE LA EXCLUSIÓN\\[2pt]
     \normalfont La deuda sobrevive a la liquidación};
  \node[favorable, below right=12mm and 4mm of excl.south] (rechaza)
    {EL TRIBUNAL RECHAZA LA EXCLUSIÓN\\[2pt]
     \normalfont La deuda ingresa al concurso y se sujeta al \art{255}};

  \node[nodo, text width=40mm, below=9mm of acoge] (ape)
    {Recurso de apelación y, en su caso, acción de protección};
  \node[nodo, text width=40mm, below=9mm of rechaza] (cas)
    {Recurso de casación del acreedor o de la \tgr{}};

  \draw[flecha] (sol) -- (excl);
  \draw[flecha] (excl.south) -- (split) -| (acoge.north);
  \draw[flecha] (excl.south) -- (split) -| (rechaza.north);
  \draw[flecha] (acoge) -- (ape);
  \draw[flecha] (rechaza) -- (cas);
\end{flujograma}

\subsection{Estrategias de defensa del deudor}

La defensa del deudor con deuda educativa depende de la sala que conozca y del estado de la
línea jurisprudencial. Frente a la postura de exclusión de la Sala Civil, las vías que la
práctica ha ensayado son: (i)~invocar la universalidad subjetiva y la regla de la \lat{lex
posterior} para sostener que la \LIR{} extingue la totalidad del pasivo; (ii)~sostener que la
especialidad del CAE sólo favorece al fisco ---no al banco--- por lo que el acreedor bancario
carece de legitimación para invocarla; y (iii)~el recurso de protección ante la Corte de
Apelaciones cuando la institución financiera mantiene registros negativos pese al término del
procedimiento. Cada una de estas vías se apoya en la evolución jurisprudencial que se
sistematiza a continuación.

%==============================================================================
\bloqueJurisprudencial
\label{sec:cap10-juris}
%==============================================================================

\subsection{Espejo comparado: cómo resuelven otros ordenamientos}

Antes de examinar la jurisprudencia chilena conviene fijar el punto de comparación, porque el
conflicto del CAE no es exclusivo de Chile: todo sistema de \discharge{} debe decidir la suerte
del crédito educativo.

\subsubsection{El modelo estadounidense: no exonerable salvo \emph{undue hardship}}

El \emph{Bankruptcy Code} \parencite{us-bankruptcy-code} resuelve la cuestión en la ley, no en
los tribunales: su \destacar{sección 523(a)(8)} declara los préstamos educativos garantizados
por entidades públicas o sin fines de lucro \destacar{no exonerables} de forma automática. El
deudor sólo puede descargarlos si inicia un \emph{adversary proceeding} y acredita que el pago
le impondría un «sufrimiento desmedido» (\emph{undue hardship}).

El estándar lo fija el célebre \falloc{Brunner v. New York State Higher Education Services
Corp.}{}\ \parencite{us-brunner}, que exige acreditar copulativamente tres elementos: (i)~que
el deudor no puede mantener un nivel de vida mínimo si paga la deuda; (ii)~que esa situación
se prolongará durante buena parte del período de amortización; y (iii)~que el deudor hizo
esfuerzos de buena fe para pagar antes de la quiebra. La \cs{} estadounidense examinó la
aplicación del test en \falloc{United Student Aid Funds, Inc. v. Espinosa}{}
\parencite{us-espinosa}.

\subsubsection{El modelo español: exoneración con protección del crédito público}

Tras la Ley 16/2022 española ---que transpone la Directiva (UE) 2019/1023---
\parencite{esp-ley16-2022, ue-directiva-2019-1023}, España sustituyó el antiguo beneficio de
exoneración del pasivo insatisfecho por un régimen de \destacar{exoneración del pasivo
insatisfecho (EPI)} con dos vías: exoneración inmediata por liquidación de la masa activa, o
sujeción a un plan de pagos de tres a cinco años que preserva la vivienda habitual.

A diferencia de Chile, el legislador español \destacar{protege expresamente el crédito
público}: la exoneración de deudas tributarias y de seguridad social tiene un tope de
10.000 euros por cada concepto, y el resto debe satisfacerse preferentemente.

\begin{foco}[La lección comparada para el debate chileno]
Estados Unidos y España resolvieron en la \emph{ley} lo que Chile dejó a los \emph{tribunales}.
Ambos excluyen ---total o parcialmente--- el crédito público o educativo del \discharge{}, pero
lo hacen con una norma expresa y con un mecanismo (el \emph{undue hardship}, el tope de
10.000 euros) que da certeza. La inestabilidad chilena no deriva de la solución de fondo
---que tiende también a la exclusión--- sino de la ausencia de una regla legal que la fije.
\end{foco}

\subsection{Evolución cronológica de la jurisprudencia de la Corte Suprema}

\subsubsection{La postura clásica de la Primera Sala (Civil): la exclusión}

La Primera Sala ha sostenido de forma persistente la \destacar{exclusión} del crédito educativo,
sobre la base del principio de especialidad normativa de los \artcc{4} y \artcc{13}: leyes
especiales como la \Lcae{} \citeyearpar{ley-20027} y la \Lfondosolidario{}
\citeyearpar{ley-19287} constituirían estatutos particulares de recuperación de fondos públicos
que desplazan las reglas de exoneración de la \LIR{}.

La línea se consolidó con \falloc{Jamarne con Salazar}{4656-2017}, de \destacar{9 de mayo de
2017} \parencite{juris-jamarne-salazar}: la Corte casó la sentencia de la \ca{Temuco} (Rol Civil
545-2016), que había confirmado el criterio de inclusión del Primer Juzgado de Letras de Temuco
(Rol C-902-2016) \parencite{juris-jamarne-primera-instancia}. El fundamento es que la \Lcae{}
constituye un \destacar{estatuto especial, orgánico, cerrado y preferente}, con mecanismos
propios para el sobreendeudamiento, la cesantía y la morosidad, y con la declaración expresa de
que las obligaciones \destacar{no prescriben} y su cobro subsiste hasta la total extinción de la
deuda. Por los \artcc{4} y \artcc{13}, la ley especial desplaza a la ley general concursal.

\subsubsection{El giro de la Tercera Sala (Constitucional): la inclusión}

Por la vía del recurso de protección, la Tercera Sala manifestó el criterio opuesto en
\falloc{Mancilla con Tesorería General de la República}{59.567-2020}, de \destacar{20 de julio
de 2020} \parencite{juris-mancilla-tgr}, originado en la apelación del recurso Rol 147-2020 de la
\ca{Puerto Montt}. La Sala Constitucional amparó al deudor: como la \LIR{} no contempla ninguna
norma expresa que exceptúe los préstamos educativos del \discharge{} ---en claro contraste con
la sección 523 del \emph{Bankruptcy Code} estadounidense, que sí los enumera---, resulta ilegal
y arbitrario que la \tgr{} o los bancos ejecuten cobros paralelos u obstaculicen la extinción
total del pasivo del deudor rehabilitado (\art{255}).

\subsubsection{La postura intermedia de las Cortes de Apelaciones}

La \ca{Concepción}, en fallos reiterados de su \destacar{Duodécima Sala}, ha rechazado recursos de
protección de estudiantes contra la suspensión del beneficio del CAE tras acogerse el deudor a una
liquidación voluntaria, sosteniendo la exclusión de la deuda regida por la \Lcae{} del mecanismo de
descarga del \art{255} por el principio de especialidad.\verificar{Individualizar roles y fechas
exactos de los fallos de la Duodécima Sala de la CA de Concepción; cotejar considerandos contra los
fallos oficiales (dato proveniente de compilación secundaria).}

\subsubsection{El retorno a la exclusión (2024-2025)}

La Primera Sala reafirmó la exclusión en pronunciamientos recientes: el \destacar{fallo de
reemplazo de 20 de diciembre de 2025} en la causa \rol{47.641-2024}, que revocó el criterio
aprobatorio de la \ca{Santiago} (Rol N.\textsuperscript{o} 9569-2024)
\parencite{juris-cs-47641-2024-reemplazo}, y la causa \rol{3557-2025}
\parencite{juris-cs-3557-2025}, sosteniendo que las obligaciones estudiantiles no se extinguen
por la resolución de término de la liquidación voluntaria. Para la sistematización de esta
línea, véanse \textcite{doc-cae-jurisprudencia} y \textcite{doc-cae-castro}.

\subsubsection{El voto disidente del Ministro Silva Cancino}

Reviste especial interés el voto disidente del Ministro \destacar{Mauricio Silva Cancino},
conforme al cual la preferencia de la \Lcae{} sólo procede cuando es invocada por el
\destacar{sujeto activo del estatuto especial ---el Estado, a través de la \tgr{}, conforme a
los artículos 12 y 13 de esa ley---} y no por las entidades bancarias privadas en su propio
interés comercial.

\begin{foco}[La tesis de la legitimación restringida como línea de defensa]
El voto de Silva Cancino ofrece al deudor una defensa concreta: cuando quien pide la exclusión
del CAE es el \emph{banco} ---y no la \tgr{}---, cabe alegar que el acreedor privado carece de
legitimación para invocar una especialidad establecida en favor del fisco. Es el argumento de
mayor recorrido práctico mientras la Primera Sala mantenga la exclusión como regla general.
\end{foco}

\subsection{La crítica dogmática a la doctrina de la exclusión}

La línea de la Sala Civil ha sido objeto de severo escrutinio académico. \textcite{caballero-cae}
subraya las implicancias del desplazamiento de la \destacar{universalidad subjetiva} en perjuicio
del deudor, y \textcite{alarcon-mala-fe} formula un análisis crítico sobre cómo la exclusión
\destacar{desarticula el interés de rehabilitación financiera} de la persona natural: perpetuar
la deuda estudiantil frente a un deudor ya liquidado contradice la finalidad misma del
\discharge{} y condena a una porción significativa de los egresados de la educación superior al
sobreendeudamiento crónico.

\begin{foco}[El nudo de fondo: la ausencia de un catálogo legal]
Toda la controversia ---dos salas de la misma Corte en sentido opuesto--- se explica por la
carencia que señala la sección \ref{sec:cap10-juris} al inicio: Chile no tiene una lista legal de
deudas excluidas del \discharge{}. Donde el legislador calló, el poder judicial actúa como
legislador de facto, y el resultado es la inseguridad jurídica de resoluciones contradictorias.
La solución de fondo ---la que proponen tanto la doctrina crítica como el derecho comparado--- no
es que los tribunales elijan, sino que la ley codifique qué deudas sobreviven y cuáles no.
\end{foco}

\begin{alerta}[Estado de la cuestión]
La materia es de criterio jurisprudencial cambiante y con sentencias muy recientes. Toda
afirmación de este capítulo debe contrastarse con la jurisprudencia vigente a la fecha de
cierre de la edición.
\end{alerta}
